\documentclass[UTF8,12pt]{ctexart}
\usepackage[a4paper,margin=24mm]{geometry}
\usepackage{amsmath,amssymb,bm,graphicx,adjustbox}
\newcommand{\fitmath}[1]{\begin{adjustbox}{max width=\linewidth}$\displaystyle #1$\end{adjustbox}}
\setlength{\parindent}{0pt}
\setlength{\parskip}{.7em}
\sloppy
\allowdisplaybreaks
\begin{document}
\section*{2005 年考研数学三 · 真题}
\subsection*{2005年全国硕士研究生招生考试试题}

\subsection*{一、填空题（本题共6小题，每小题4分，满分24分）}

（1）极限 \(\displaystyle \lim_{x\to\infty}x\sin\dfrac{\displaystyle 2x}{\displaystyle x^2+1}=\)\_\_\_\_\_\_。


（2）微分方程 \(\displaystyle xy^{\prime}+y=0\) 满足初始条件 \(\displaystyle y(1)=2\) 的特解为\_\_\_\_\_\_。


（3）设二元函数 \(\displaystyle z=xe^{x+y}+(x+1)\ln(1+y)\)，则 \(\displaystyle \left.\mathrm dz\right|_{(1,\allowbreak 0)}=\)\_\_\_\_\_\_。


（4）设行向量组 \(\displaystyle (2,\allowbreak 1,\allowbreak 1,\allowbreak 1),\allowbreak (2,\allowbreak 1,\allowbreak a,\allowbreak a),\allowbreak (3,\allowbreak 2,\allowbreak 1,\allowbreak a),\allowbreak (4,\allowbreak 3,\allowbreak 2,\allowbreak 1)\) 线性相关，且 \(\displaystyle a\ne1\)，则 \(\displaystyle a=\)\_\_\_\_\_\_。


（5）从数 \(\displaystyle 1,\allowbreak 2,\allowbreak 3,\allowbreak 4\) 中任取一个数，记为 \(\displaystyle X\)，再从 \(\displaystyle 1,\allowbreak \ldots,\allowbreak X\) 中任取一个数，记为 \(\displaystyle Y\)，则 \(\displaystyle P\{Y=2\}=\)\_\_\_\_\_\_。


（6）设二维随机变量 \(\displaystyle (X,\allowbreak Y)\) 的概率分布为 \(\displaystyle \begin{array}{c|cc}X\backslash Y&0&1\\\hline0&0.4&a\\1&b&0.1\end{array}\)。若随机事件 \(\displaystyle \{X=0\}\) 与 \(\displaystyle \{X+Y=1\}\) 相互独立，则 \(\displaystyle a=\)\_\_\_\_\_\_，\(\displaystyle b=\)\_\_\_\_\_\_。


\subsection*{二、选择题（本题共8小题，每小题4分，满分32分）}

（7）当 \(\displaystyle a\) 取下列哪个值时，函数 \(\displaystyle f(x)=2x^3-9x^2+12x-a\) 恰有两个不同的零点（　）。（A）2。（B）4。（C）6。（D）8。


（8）设 \(\displaystyle I_1=\displaystyle \iint_D\cos\sqrt{x^2+y^2}\,\mathrm d\sigma\)，\(\displaystyle I_2=\displaystyle \iint_D\cos(x^2+y^2)\,\mathrm d\sigma\)，\(\displaystyle I_3=\displaystyle \iint_D\cos(x^2+y^2)^2\,\mathrm d\sigma\)，其中 \(\displaystyle D=\{(x,\allowbreak y)\mid x^2+y^2\le1\}\)，则（　）。（A）\(\displaystyle I_3>I_2>I_1\)。（B）\(\displaystyle I_1>I_2>I_3\)。（C）\(\displaystyle I_2>I_1>I_3\)。（D）\(\displaystyle I_3>I_1>I_2\)。


（9）设 \(\displaystyle a_n>0\)，\(\displaystyle n=1,\allowbreak 2,\allowbreak \ldots\)，若 \(\displaystyle \sum_{n=1}^{\infty}a_n\) 发散，\(\displaystyle \sum_{n=1}^{\infty}(-1)^{n-1}a_n\) 收敛，则下列结论正确的是（　）。（A）\(\displaystyle \sum a_{2n-1}\) 收敛，\(\displaystyle \sum a_{2n}\) 发散。（B）\(\displaystyle \sum a_{2n}\) 收敛，\(\displaystyle \sum a_{2n-1}\) 发散。（C）\(\displaystyle \sum(a_{2n-1}+a_{2n})\) 收敛。（D）\(\displaystyle \sum(a_{2n-1}-a_{2n})\) 收敛。


（10）设 \(\displaystyle f(x)=x\sin x+\cos x\)，下列命题中正确的是（　）。（A）\(\displaystyle f(0)\) 是极大值，\(\displaystyle f(\dfrac{\displaystyle \pi}{\displaystyle 2})\) 是极小值。（B）\(\displaystyle f(0)\) 是极小值，\(\displaystyle f(\dfrac{\displaystyle \pi}{\displaystyle 2})\) 是极大值。（C）两者都是极大值。（D）两者都是极小值。


（11）以下四个命题中，正确的是（　）。


（11）（续）（A）若 \(\displaystyle f^{\prime}(x)\) 在 \(\displaystyle (0,\allowbreak 1)\) 内连续，则 \(\displaystyle f(x)\) 在 \(\displaystyle (0,\allowbreak 1)\) 内有界。（B）若 \(\displaystyle f(x)\) 在 \(\displaystyle (0,\allowbreak 1)\) 内连续，则 \(\displaystyle f(x)\) 在 \(\displaystyle (0,\allowbreak 1)\) 内有界。（C）若 \(\displaystyle f^{\prime}(x)\) 在 \(\displaystyle (0,\allowbreak 1)\) 内有界，则 \(\displaystyle f(x)\) 在 \(\displaystyle (0,\allowbreak 1)\) 内有界。（D）若 \(\displaystyle f(x)\) 在 \(\displaystyle (0,\allowbreak 1)\) 内有界，则 \(\displaystyle f^{\prime}(x)\) 在 \(\displaystyle (0,\allowbreak 1)\) 内有界。


（12）设矩阵 \(\displaystyle A=(a_{ij})_{3\times3}\)，满足 \(\displaystyle A^*=A^{\mathrm T}\)，其中 \(\displaystyle A^*\) 为 \(\displaystyle A\) 的伴随矩阵，\(\displaystyle A^{\mathrm T}\) 为 \(\displaystyle A\) 的转置矩阵。若 \(\displaystyle a_{11},\allowbreak a_{12},\allowbreak a_{13}\) 为三个相等的正数，则 \(\displaystyle a_{11}\) 为（　）。（A）\(\displaystyle \dfrac{\displaystyle \sqrt3}{\displaystyle 3}\)。（B）\(\displaystyle 3\)。（C）\(\displaystyle \dfrac{\displaystyle 1}{\displaystyle 3}\)。（D）\(\displaystyle \sqrt3\)。


（13）设 \(\displaystyle \lambda_1,\allowbreak \lambda_2\) 是矩阵 \(\displaystyle A\) 的两个不同的特征值，对应的特征向量分别为 \(\displaystyle \alpha_1,\allowbreak \alpha_2\)，则 \(\displaystyle \alpha_1,\allowbreak A(\alpha_1+\alpha_2)\) 线性无关的充分必要条件是（　）。（A）\(\displaystyle \lambda_1=0\)。（B）\(\displaystyle \lambda_2=0\)。（C）\(\displaystyle \lambda_1\ne0\)。（D）\(\displaystyle \lambda_2\ne0\)。


（14）（超纲题）设一批零件的长度服从正态分布 \(\displaystyle N(\mu,\allowbreak \sigma^2)\)，其中 \(\displaystyle \mu,\allowbreak \sigma^2\) 均未知。现从中随机抽取16个零件，测得样本均值 \(\displaystyle \bar x=20\)（cm），样本标准差 \(\displaystyle S=1\)（cm），则 \(\displaystyle \mu\) 的置信度为0.90的置信区间是（　）。（A）\(\displaystyle (20-\dfrac{\displaystyle 1}{\displaystyle 4}t_{0.05}(16),\allowbreak 20+\dfrac{\displaystyle 1}{\displaystyle 4}t_{0.05}(16))\)。（B）\(\displaystyle (20-\dfrac{\displaystyle 1}{\displaystyle 4}t_{0.1}(16),\allowbreak 20+\dfrac{\displaystyle 1}{\displaystyle 4}t_{0.1}(16))\)。（C）\(\displaystyle (20-\dfrac{\displaystyle 1}{\displaystyle 4}t_{0.05}(15),\allowbreak 20+\dfrac{\displaystyle 1}{\displaystyle 4}t_{0.05}(15))\)。（D）\(\displaystyle (20-\dfrac{\displaystyle 1}{\displaystyle 4}t_{0.1}(15),\allowbreak 20+\dfrac{\displaystyle 1}{\displaystyle 4}t_{0.1}(15))\)。


\subsection*{三、解答题（本题共9小题，满分94分。解答应写出文字说明、证明过程或演算步骤）}

（15）（本题满分8分）求 \(\displaystyle \lim_{x\to0}\left(\dfrac{\displaystyle 1+x}{\displaystyle 1-e^{-x}}-\dfrac{\displaystyle 1}{\displaystyle x}\right)\)。


（16）（本题满分8分）设 \(\displaystyle f(u)\) 具有二阶连续导数，且 \(\displaystyle g(x,\allowbreak y)=f(\dfrac{\displaystyle y}{\displaystyle x})+yf(\dfrac{\displaystyle x}{\displaystyle y})\)，求 \(\displaystyle x^2\dfrac{\displaystyle \partial^2g}{\displaystyle \partial x^2}-y^2\dfrac{\displaystyle \partial^2g}{\displaystyle \partial y^2}\)。


（17）（本题满分9分）计算二重积分 \(\displaystyle \iint_D|x^2+y^2-1|\,\mathrm d\sigma\)，其中 \(\displaystyle D=\{(x,\allowbreak y)\mid0\le x\le1,\allowbreak 0\le y\le1\}\)。


（18）（本题满分9分）求幂级数 \(\displaystyle \sum_{n=1}^{\infty}(\dfrac{\displaystyle 1}{\displaystyle 2n+1}-1)x^{2n}\) 在区间 \(\displaystyle (-1,\allowbreak 1)\) 内的和函数 \(\displaystyle S(x)\)。


（19）（本题满分8分）设 \(\displaystyle f(x),\allowbreak g(x)\) 在 \(\displaystyle [0,\allowbreak 1]\) 上的导数连续，且 \(\displaystyle f(0)=0\)，\(\displaystyle f^{\prime}(x)\ge0\)，\(\displaystyle g^{\prime}(x)\ge0\)。证明：对任何 \(\displaystyle a\in[0,\allowbreak 1]\)，有 \(\displaystyle \int_0^ag(x)f^{\prime}(x)\,\mathrm dx+\displaystyle \int_0^1f(x)g^{\prime}(x)\,\mathrm dx\ge f(a)g(1)\)。


（20）（本题满分13分）已知齐次线性方程组


（20）（续）（i）\(\displaystyle \begin{cases}x_1+2x_2+3x_3=0,\allowbreak \\2x_1+3x_2+5x_3=0,\allowbreak \\x_1+x_2+ax_3=0,\allowbreak \end{cases}\) 和（ii）\(\displaystyle \begin{cases}x_1+bx_2+cx_3=0,\allowbreak \\2x_1+b^2x_2+(c+1)x_3=0.\end{cases}\) 同解，求 \(\displaystyle a,\allowbreak b,\allowbreak c\) 的值。


（21）（本题满分13分）设 \(\displaystyle D=\begin{pmatrix}A&C\\C^{\mathrm T}&B\end{pmatrix}\) 为正定矩阵，其中 \(\displaystyle A,\allowbreak B\) 分别为 \(\displaystyle m\) 阶、\(\displaystyle n\) 阶对称矩阵，\(\displaystyle C\) 为 \(\displaystyle m\times n\) 矩阵。（I）计算 \(\displaystyle P^{\mathrm T}DP\)，其中 \(\displaystyle P=\begin{pmatrix}E_m&-A^{-1}C\\O&E_n\end{pmatrix}\)；（II）利用（I）的结果判断矩阵 \(\displaystyle B-C^{\mathrm T}A^{-1}C\) 是否为正定矩阵，并证明你的结论。


（22）（本题满分13分）设二维随机变量 \(\displaystyle (X,\allowbreak Y)\) 的概率密度为 \(\displaystyle f(x,\allowbreak y)=\begin{cases}1,\allowbreak &0<x<1,\allowbreak \ 0<y<2x,\allowbreak \\0,\allowbreak &\text{其他}.\end{cases}\) 求：（I）\(\displaystyle (X,\allowbreak Y)\) 的边缘概率密度 \(\displaystyle f_X(x),\allowbreak f_Y(y)\)；（II）\(\displaystyle Z=2X-Y\) 的概率密度 \(\displaystyle f_Z(z)\)；（III）\(\displaystyle P\{Y\le\dfrac{\displaystyle 1}{\displaystyle 2}\mid X\le\dfrac{\displaystyle 1}{\displaystyle 2}\}\)。


（23）（本题满分13分）设 \(\displaystyle X_1,\allowbreak X_2,\allowbreak \ldots,\allowbreak X_n\;(n>2)\) 为来自总体 \(\displaystyle N(0,\allowbreak \sigma^2)\) 的简单随机样本，其样本均值为 \(\displaystyle \bar X\)。记 \(\displaystyle Y_i=X_i-\bar X\)，\(\displaystyle i=1,\allowbreak 2,\allowbreak \ldots,\allowbreak n\)。（I）求 \(\displaystyle Y_i\) 的方差 \(\displaystyle D(Y_i)\)；（II）求 \(\displaystyle Y_1\) 与 \(\displaystyle Y_n\) 的协方差 \(\displaystyle \operatorname{Cov}(Y_1,\allowbreak Y_n)\)；（III）\(\displaystyle c(Y_1+Y_n)^2\) 是 \(\displaystyle \sigma^2\) 的无偏估计量，求常数 \(\displaystyle C\)。（无偏估计为超纲概念，可改为“若 \(\displaystyle E(c(Y_1+Y_n)^2)=\sigma^2\)，求常数 \(\displaystyle C\)”）

\end{document}
